\documentclass[11pt,a4paper]{article}

\usepackage[utf8]{inputenc}
\usepackage[T1]{fontenc}
\usepackage{amsmath,amssymb}
\usepackage{graphicx}
\usepackage{booktabs}
\usepackage[margin=2.5cm]{geometry}
\usepackage[hidelinks]{hyperref}

\title{UNET Inference \& Prediction Pipeline \\ {\normalsize AzONetV2 — Thickness Estimation from Spectra}}
\author{}
\date{\today}

\begin{document}
\maketitle

\begin{abstract}
This report describes the inference pipeline in \texttt{notebooks/6\_results\_colab/Predictions.py}, which loads a pre-trained UNET model to predict film thicknesses from both simulated and experimental spectral data. Predictions are generated for the full train/test/validation splits as well as external experimental 145\,nm spectra, with all results exported as NumPy arrays.
\end{abstract}

%---------------------------------------------------------------------
\section{Introduction}

The AzONetV2 project employs deep learning to estimate thin-film thickness from reflectance/transmission spectra. The script \texttt{Predictions.py} is the final stage of the pipeline: it loads the trained UNET model and runs inference on all available datasets, producing predictions that are subsequently used for error analysis and model evaluation.

%---------------------------------------------------------------------
\section{Data}

\subsection{Simulated Datasets}

Train, test, and validation splits are stored as HDF5 files under:

\begin{verbatim}
notebooks/3_data_simulation/NewDataAzONetV2/Sets/
\end{verbatim}

Each \texttt{.h5} file contains two datasets: \texttt{x\_total} (spectra) and \texttt{y\_total} (thickness labels). The function \texttt{load\_data\_normalization\_sample\_General} loads these splits and applies per-sample min-max normalization via \texttt{normalization\_by\_sample}.

\subsection{Experimental Data}

Experimental spectra for 145\,nm thin films are loaded from:

\begin{verbatim}
results/dataframe_spectrum_thickness_145_final.pkl
\end{verbatim}

Ground-truth thickness values for these samples come from differential evolution parameters stored in:

\begin{verbatim}
results/SciPy_IIM/differential_evolution_NonLinear_145F.npy
\end{verbatim}

Both datasets undergo the same per-sample normalization before inference.

%---------------------------------------------------------------------
\section{Methodology}

\subsection{Per-Sample Normalization}

The function \texttt{normalization\_by\_sample} applies min-max scaling independently to each sample:

\[
x_{\text{norm}} = \frac{x - \min(x)}{\max(x) - \min(x) + \epsilon}, \quad \epsilon = 10^{-6}
\]

This ensures each spectrum is normalized to the $[0, 1]$ range without cross-sample information leakage.

\subsection{UNET Model}

The pre-trained UNET model is loaded from:

\begin{verbatim}
notebooks/6_results_colab/models/Copia de model.keras
\end{verbatim}

An integrity check is performed by evaluating the model on the test set using \texttt{modelo\_UNET.evaluate}. The model architecture follows a standard encoder-decoder UNET design suitable for 1D spectral regression.

\subsection{GPU Allocation}

GPU resources are configured via TensorFlow's \texttt{set\_visible\_devices}, targeting the first available GPU device. This is essential for efficient batch inference on the full datasets.

%---------------------------------------------------------------------
\section{Predictions \& Outputs}

Predictions are generated using \texttt{modelo\_UNET.predict} for all four data subsets:

\begin{itemize}
    \item \textbf{Train}: predictions on the full training split
    \item \textbf{Test}: predictions on the held-out test split
    \item \textbf{Validation}: predictions on the validation split
    \item \textbf{Experimental 145\,nm}: predictions on external experimental spectra
\end{itemize}

All outputs are saved under the \texttt{MSE\_data/} directory (the prefix \texttt{MSE} is set by the \texttt{\_name} variable at the top of the script):

\begin{table}[h]
\centering
\begin{tabular}{ll}
\toprule
\textbf{File} & \textbf{Description} \\
\midrule
\texttt{preds\_train.npy} & Train set predictions \\
\texttt{y\_train.npy}     & Train set ground truth \\
\texttt{preds\_test.npy}  & Test set predictions \\
\texttt{y\_test.npy}      & Test set ground truth \\
\texttt{preds\_val.npy}   & Validation set predictions \\
\texttt{y\_val.npy}       & Validation set ground truth \\
\texttt{preds\_145.npy}   & Experimental predictions \\
\texttt{y\_145.npy}       & Experimental ground truth \\
\bottomrule
\end{tabular}
\caption{Saved output files under \texttt{MSE\_data/}.}
\label{tab:outputs}
\end{table}

%---------------------------------------------------------------------
\section{Reproducibility}

\begin{itemize}
    \item \textbf{Environment}: Google Colab with GPU runtime
    \item \textbf{TensorFlow version}: as available in the Colab environment at runtime
    \item \textbf{Model file}: \texttt{Copia de model.keras} (pre-trained UNET)
    \item \textbf{Seed}: no explicit random seed is set; results are deterministic given the same model weights and data
\end{itemize}

To reproduce the results, mount Google Drive, verify all input paths in the script, and execute \texttt{Predictions.py}. The output \texttt{.npy} files will be written to the \texttt{MSE\_data/} folder within the same notebook directory.

\end{document}
